% !TEX root = ../../main-es.tex
% chapters/es/04-teorema.tex
\chapter{No-interferencia composicional y soundness del verificador}
\criterio{plan}{15}\criterio{analisis}{20}
\label{cap:teorema}

\porque{Rigor metodológico + análisis (EICT 15\% + 20\%).}{%
  El corazón formal de la tesis. Fallo típico: enunciar un teorema fuerte cuya prueba es
  trivial o, peor, falsa. La estrategia es enunciar lo \emph{mínimo} que el código garantiza
  (no-interferencia por aislamiento de stores) y demostrar eso limpiamente.}

\section{Definiciones previas}
\porque{Setup.}{Store, snapshot, frontera de celda, el fragmento fork-join $F_{\mathrm{fj}}$.}

\section{Teorema de No-Interferencia Composicional}
\begin{teorema}[No-interferencia en $F_{\mathrm{fj}}$]\label{thm:noi}
  Para todo $f,g,h,k$ en el fragmento fork-join con stores aislados,
  \[
    (f \tensor g) \seqcomp (h \tensor k) \;=\; (f \seqcomp h) \tensor (g \seqcomp k).
  \]
\end{teorema}
\porque{Enunciado.}{La ley de intercambio. Caso finito del monoide simétrico.}

\section{Lemas}
\porque{Descomposición.}{Aislamiento $\Rightarrow$ escrituras de una rama no son visibles para la otra; el join preserva el orden de la lista de targets.}

\section{Demostración}
\porque{Prueba por inducción.}{Sobre la estructura de los morfismos. Casos: identidad, CALL secuencial, fork, join. Recomendación: mecanización Lean del caso finito como anexo opcional de alto valor.}

\section{Corolario: inalcanzabilidad de una clase de fallos}
\begin{corolario}\label{cor:contaminacion}
  La clase de fallos por contaminación de contexto entre ramas fork es vacía en
  $F_{\mathrm{fj}}$.
\end{corolario}
\porque{Del teorema al claim medible.}{Es la primera fila de la taxonomía de O4 que pasa a ``inalcanzable por construcción''.}

\section{Soundness del verificador estático (O3)}
\porque{Segundo claim formal.}{Si el verificador rechaza \cellref{C}{f} entonces, bajo la semántica de §\ref{sec:semantica}, toda ejecución que alcance esa referencia fallaría. Prueba por inducción en la derivación de tipado. Completitud acotada por el fragmento estáticamente resoluble.}
